当 \(p=2\) 时，Plancherel 定理和 \(|I|=1\) 给出结论，此时 \(\varepsilon=0\)。设 \(2<p<\infty\) 且 \(s=s_+(p)+\varepsilon\)。令 \(g=J^s f\)。在频率空间中取一个光滑单位分解，由紧支撑的低频截断 \(\chi_0\) 与二进环带截断 \(\psi(2^{-j}\xi)\) 组成，其中 \(j\ge1\)，并使用同一个固定的光滑紧支撑环带符号 \(\psi\)。选取光滑紧支撑的环带符号 \(\widetilde\psi\)，使它在 \(\operatorname{supp}\psi\) 上等于一。
乘子
\[
 m_j(\xi)=2^{js}\widetilde\psi(2^{-j}\xi)(1+|\xi|^2)^{-s/2}
\]
具有一致可积的核。事实上，尺度变换后，
\(m_j(2^j\eta)=\widetilde\psi(\eta)(2^{-2j}+|\eta|^2)^{-s/2}\)
的支撑落在固定的且远离零点的紧集内，各阶导数一致有界。因此反复分部积分给出逆变换的一致 \(L^1\) 界。尺度 \(2^j\) 的高频波为 \(2^{-js}U(T)\psi(2^{-j}D)m_j(D)g\)。
其时空范数不超过
\[
 C\,2^{-j\varepsilon}(1+j)^{\beta(p)}\|g\|_{L^p(\mathbb R^n)},
 \qquad
 \beta(p)=\min\{(n-1)(1-2/p),2/p\},
\]
这里使用了 \ref{ms:annular-lp} 和 Young 不等式。这些界对 \(j\ge1\) 可求和。

对于低频波，其乘子
\(e^{-iT|\xi|}\chi_0(\xi)(1+|\xi|^2)^{-s/2}\)
一致满足 Mihlin 导数界，其中 \(T\in I\)：每个 \(|\xi|^{|\alpha|}|\partial_\xi^\alpha m(\xi)|\) 在 \(\xi\ne0\) 上都一致有界，直至 \(\lfloor n/2\rfloor+1\) 阶。因此，命题~\ref{ext:mihlin-multiplier} 给出一致的空间 \(L^p\) 界。由于 \(|I|=1\)，此界也控制时空范数。将低频和高频的界相加，证得该估计。对 Schwartz 数据，频率分解在分布意义下收敛到 \(U(T)f\)，而由范数上界可求和可知，\(L^p\) 极限与这一分布相同。
